% Part: first-order-logic
% Chapter: lindstrom
% Section: ls-property

\documentclass[../../../include/open-logic-section]{subfiles}

\begin{document}

\olfileid[pt]{mod}{lin}{lsp}

\olsection{Propriedades da compacidade e de L\"owenheim--Skolem}

Apresentamos agora as extensões naturais da compacidade e de
L\"owenheim--Skolem ao caso das lógicas abstratas.

\begin{defn}
Uma lógica abstrata $\tuple{L, \models_L}$ tem a \emph{propriedade da
  compacidade} se todo conjunto $\Gamma$ de
$L(\Lang{L})$-!!{sentence}s é satisfatível sempre que todo
$\Gamma_0 \subseteq \Gamma$ finito é satisfatível.
\end{defn}

\begin{defn}
$\tuple{L, \models_L}$ tem a \emph{propriedade descendente de
  L\"owenheim--Skolem} se todo $\Gamma$ satisfatível tem um modelo
!!{enumerable}.
\end{defn}


A noção de isomorfismo parcial de
\olref[bas][pis]{defn:partialisom} é puramente ``algébrica'' (isto
é, formulada sem referência às !!{sentence}s da linguagem, apenas às
constantes fornecidas pela !!{language}~$\Lang{L}$ das
!!{structure}s) e, portanto, aplica-se ao caso das lógicas abstratas.
No caso da lógica de primeira ordem, sabemos por
\olref[bas][pis]{thm:p-isom2} que, se duas !!{structure}s são
parcialmente isomorfas, então são elementarmente equivalentes. Essa
demonstração não se estende às lógicas abstratas, pois a indução em
!!{formula}s pode não estar disponível para um
$!E \in L(\Lang{L})$ arbitrário; ainda assim, o teorema é verdadeiro,
desde que valha a propriedade de L\"owenheim--Skolem.

\begin{thm}
\ollabel{thm:abstract-p-isom}
Suponha que $\tuple{L, \models_L}$ seja uma lógica normal com a
propriedade de L\"owenheim--Skolem. Então quaisquer duas !!{structure}s
parcialmente isomorfas são elementarmente equivalentes em
$\tuple{L, \models_L}$.
\end{thm}

\begin{proof}
Suponha que $\Struct{M} \simeq_p \Struct{N}$, mas, para algum $!E$,
$\Struct{M} \models_L !E$ enquanto $\Struct{N} \not\models_L !E$. Pela
Propriedade do Isomorfismo, podemos supor que $\Domain{M}$ e
$\Domain{N}$ sejam disjuntos; pela Propriedade da Expansão, podemos
supor que $!E \in L(\Lang{L})$ para uma !!{language} finita
$\Lang{L}$. Seja $\mathcal{I}$ um conjunto de isomorfismos parciais
entre $\Struct{M}$ e $\Struct{N}$ e, sem perda de generalidade, suponha
também que, se $p \in \mathcal{I}$ e $q \subseteq p$, então
$q \in \mathcal{I}$.

$\Domain{M}^{<\omega}$ é o conjunto das sequências finitas de
!!{element}s de~$\Domain{M}$. Seja $S$ a relação ternária sobre
$\Domain{M}^{<\omega}$ que representa a concatenação, isto é, se
$\mathbf{a}, \mathbf{b}, \mathbf{c} \in \Domain{M}^{<\omega}$, então
$S(\mathbf{a}, \mathbf{b}, \mathbf{c})$ vale se, e somente se,
$\mathbf{c}$ é a concatenação de $\mathbf{a}$ e $\mathbf{b}$; seja $T$
a relação ternária tal que $T(\mathbf{a}, b, \mathbf{c})$ vale para
$b \in M$ e $\mathbf{a}, \mathbf{c} \in \Domain{M}^{<\omega}$ se, e
somente se, $\mathbf{a} = a_1, \dots a_n$ e
$\mathbf{c} = a_1, \dots a_n, b$. Escolha novos !!{predicate}s
ternários $P$ e $Q$ e forme a !!{structure} $\Struct{M}^*$ com
universo $\Domain{M} \cup \Domain{M}^{<\omega}$, tendo $\Struct{M}$
como subestrutura e interpretando $P$ e $Q$ pelas relações de
concatenação $S$ e $T$ (assim, $\Struct{M}^*$ está na !!{language}
$\Lang{L} \cup \{ P, Q\}$).

Defina $\Domain{N}^{<\omega}$, $S'$, $T'$, $P'$, $Q'$ e
$\Struct{N}^*$ de modo análogo. Como, por hipótese,
$\Struct{M} \simeq_p \Struct{N}$, existe uma relação $I$ entre
$\Domain{M}^{<\omega}$ e $\Domain{N}^{<\omega}$ tal que
$I(\mathbf{a}, \mathbf{b})$ vale se, e somente se, $\mathbf{a}$ e
$\mathbf{b}$ são isomorfas e satisfazem a condição de vaivém de
\olref[bas][pis]{defn:partialisom}. Agora, seja $\Struct{M}$ a
!!{structure} cujo !!{domain} é a união dos !!{domain}s de
$\Struct{M}^*$ e $\Struct{N}^*$, tendo $\Struct{M}^*$ e $\Struct{N}^*$
como sub!!{structure}s, na !!{language} com um !!{predicate} binário
adicional~$R$, interpretado pela relação~$I$, e !!{predicate}s que
denotam os !!{domain}s $\Domain{M}^*$ e~$\Domain{N}*$.

\begin{figure}[h]
  \centering
  \begin{tikzpicture}[node distance=2cm, auto, thick, >=stealth']
    \draw [rounded corners] (0,0) -- (8,0) -- (8,4) -- (0,4) --  cycle;
    \draw (2,2) circle (0.5cm);
    \draw (2,2) circle (1.25cm);
    \draw (6,2) circle (0.5cm);
    \draw (6,2) circle (1.25cm);
    \path node at (0.75,3.5) {\large $\Struct{M}$};
    \path node at (2,2) {\large $\Struct{M}$};
    \path node at (6,2) {\large $\Struct{N}$};
    \path node at (3.5,1) {\large $\Struct{M}^*$};
    \path node at (7.5,1) {\large $\Struct{N}^*$};
    \node (Idom) at (2.8,2) {};
    \node (Irng) at (5.2,2) {};
    \draw[<->, bend left] (Idom) to node {\large $I$} (Irng) ;
  \end{tikzpicture} 
  \caption{A !!{structure}~$\Struct{M}$ com o isomorfismo parcial
    interno.}
\end{figure}

A observação crucial é que, na !!{language} da
!!{structure}~$\Struct{M}$, existe uma !!{sentence} \emph{de primeira
  ordem} $!D_1$, verdadeira em $\Struct{M}$, que afirma
$\Struct{M} \models_L !E$ e $\Struct{N} \not\models_L !E$ (isso exige
a Propriedade da Relativização), bem como uma !!{sentence}
\emph{de primeira ordem} $!D_2$, verdadeira em $\Struct{M}$, que
afirma que $\Struct{M} \simeq_p \Struct{N}$ por meio do isomorfismo
parcial~$I$. Pela Propriedade de L\"owenheim--Skolem, $!D_1$ e $!D_2$
são conjuntamente verdadeiras em um modelo !!{enumerable}
$\Struct{M}_0$ que contém subestruturas parcialmente isomorfas
$\Struct{M}_0$ e $\Struct{N}_0$, tais que
$\Struct{M}_0 \models_L !E$ e $\Struct{N}_0 \not\models_L !E$. Mas
no caso !!{enumerable}, !!{structure}s parcialmente isomorfas são, de fato,
isomorfas por \olref[bas][pis]{thm:p-isom1}, o que contradiz a
Propriedade do Isomorfismo das lógicas abstratas normais.
\end{proof}

\end{document}
